\newpage
\chapter{Позадина и поврзани дела}
Литературата за симулација на меки тела има произведено повеќе различни фамилии на методи. Тие се разликуваат во начинот на кој го дискретизираат телото, како ги поврзуваат силите со деформацијата и како интегрираат со текот на времето. Разбирањето на нивните предности и недостатоци е неопходно пред да се дискутираат колизиите, бидејќи различни техники за колизии функционираат добро за различни парадигми

\section{Парадигми за симулација на меки тела}
\subsection{Системи на маса и пружина}
Тие се најраните и наједноставните модели. Телото е дискретизирано како множество од точкести маси (point masses) кои се поврзани со Хукови пружини; силите на пружината го почитуваат Хуковиот закон и се интегрираат нанапред во времето, најчесто со експлицитна Ојлерова или Верле интеграција. Provot \cite{Provot1995DeformationCI} го пионира овојо пристап за ткаенина, со структурни, смолкнувачки и искривувачки пружини што ја формираат основата на моделот. Бараф и Виткин \cite{BaraffDavidWitkin1998} го прошируваат моделот со имплицитна интеграција за да се премина бариерата на експлицитна стабилност. 

Јачината на овие системи е во нивната едноставност: имаат едноставен математички модел и податочните структури се тривијални. Слабостите се тоа што цврстината(stiffness) заедно со експлицитна интеграција произведува нестабилност, што форсира користење на мали временски чекори. Однесувањето е анисотропично зависно од топологијата, изборот на триангулација на ткаенината има ефект на тоа како таа се однесува. Подесување на систем на маса и пружина за да соодветствува на вистински материјал е многу тешко. И покрај овие недостатоци, овој модел се уште е чест во симулацијата на ткаенина и во едукативни имплементации токму поради тоа што е едноставен.

\subsection{Метод на конечни елементи (Finite Element Method)}
Методот на конечни елементи (оттука FEM) е главниот пристап на механика на континуум. Телото се дискретизира во елементи (тетраедри во 3Д, триаголници во 2Д) и симулацијата решава за полето на изместување (displacement field) што минимизира функција на енергија што се изведува од напрегањето. Има 3 фамилии кои се важни за меки тела:

\textit{Линеарен FEM} користи линеарна(Коши) мерка за напрегање. Тој е брз бидејќи може да се изврши предпресметка на повеќето елементи од почетната конфигурација. Проблемот е што не е ротационо инваријатен овој метод, тело што се ротира а нема деформација ќе регистрира како да нема напрегање, што произведува невидливи сили што го туркаат телото во погрешни насоки.

\textit{Коротациски FEM} (Милер и Грос \cite{MullerAndGross}) го решава проблемот на ротација со екстрахирање на ротација по елемент од градиентот на деформација и аплицирање на линеарерно напрегање во ротираната рамка. Резултатот е ротационо инваријантент и стабилен за деформации од нормална голема, со средно покачување на цената за пресметка. Коротацискиот FEM е де факто стандард за FEM меки тела во реално време во интерактивни апликации.

\textit{Нелинеарен FEM} користи напрегање од Green или други видови на мерења за напрегање заедно со хипереластична густина на енергја(Нео-Хукова, Муни-Ривлин). Методот е физички точен, но скап: градиентот е покомплексен и стабилната интеграција најчесто бара имплицитни методи кои решаваат линеарен или нелинеарен систем на секој временски чекор.

\subsection{Спојување на облици (Shape matching)}
Спојувањето на облици \cite{MullerAndGross} е метод без мрежи (meshless method). Деформираната конфигурација на кластер од честички се спојува со конфигурацијата во мирување со пресметување на цврстата трансформација со најдобро прилагодување (best fit), по што честичките се повлекуваат кон трансформираните почетни позиции. Методот е ефтин, многу стабилен и елегантно се справува со големи деформации и ротации. Главното ограничување е тоа што контролира глобален облик на телото наместо локални напрегања, што прави мали одлики на телото тешко да се прикажат. NVIDIA Flex користи спојување на облици за цврсти тела.

\subsection{Динамика базирана на позиции (Position-Based Dynamics)}
\textbf{Диманика базирана на позиции (понатаму PBD)}\cite{pbd} има различен пристап: наместо да се изведуваат силите од енергијата и да се интегрира, директно проектира честички за да се задоволат геометриски ограничувања. После чекор на предвидување на позиција во Симплектичен-Ојлер стил (Symplectic-Euler), итеративен Гаус-Сајдел решавач го посетува секое ограничување(растојание, волумен, виткање, колизии) и ги корегира позициите на честичките за да се задоволи ограничувањето. Брзините на честичките се добиваат од промената на позицијата на крајот на временскиот чекор.

Резултатот е безусловно стабилен и лесен да се надгради со додатни типови на ограничувања. PBD доминира во симулациите на меки тела во реално време во игри и виртуелната реалност, делумно бидејќи бројот на итерации дава интиутивна контрола на квалитетот и делумно бидејќи формулациите базирани на ограничувања за ткаенини, јажиња, коса, течности и грануларни материјали сите може да се вклопат во еден решавач.

Главното ограничување на обичниот PBD е што ефективната цврстина зависи од бројот на итерации: удвојувањето на итерациите ефективно ја удвојува и цврстината на телата. Ова прави потешкотии при специфицирање на конкретен материјал, и корисникот мора да ги подесува итерациите и временскиот чекор заедно за да ја добие посакуваната цврстина.

\subsection{Проширена диманика базирана на позиции (Extended PBD)}
\textbf{Extended PBD (понатаму XPBD)} го реформулира PBD како ограничен оптимизациски проблем со усогласеност (compliance), при што акумулира Лагранжови множителни низ итерациите. Параметарот за усогласеност $\alpha$ означува единица на инверзна цврстина, а во ажурирањето на позиции во секоја итерација се користи $\tilde{\alpha} = \frac{\alpha}{dt^2}$, што ја прави цврстината на телото независна од бројот на итерации. 

XPBD ја задржува стабилноста на PBD и можноста за надоградување при што се додава параметар за полесно специфицирање на специфични материјали. Тој е интеграторот што се користи во овој труд и формулацијата се објаснува формално во глава 3.

\begin{table}[h]
	\centering
	\begin{tabularx}{\textwidth}{lXXXX}
		\toprule
		\textbf{Paradigm} & \textbf{Cost} & \textbf{Stability under stiffness} & \textbf{Material control} & \textbf{Typical use} \\
		\midrule
		Mass-spring & Very low & Poor with explicit integration & Coarse, topology-dependent & Cloth, education \\
		Linear FEM & Low & Moderate & Continuum strain & Small deformations \\
		Corotational FEM & Moderate & Good & Continuum strain, rotation-invariant & Real-time soft bodies, surgical sim \\
		Nonlinear FEM & High & Good (with implicit) & Hyperelastic & Offline accurate sim \\
		Shape matching & Very low & Excellent & Global, not local & Stylised solids, Flex \\
		PBD & Low & Excellent & Compliance, iteration-coupled & Real-time everything \\
		XPBD & Low & Excellent & Compliance, iteration-independent & Real-time everything (preferred) \\
		\bottomrule
	\end{tabularx}
	\caption{Comparison of simulation paradigms}
	\label{tab:comparison}
\end{table}

\section{Историја на методи базирани на позиција}
Методите базирани на позиција сега се широко прифатени во графиката во реално време, но имале инкрементален развој низ годините. Секој од овие чекори решавал проблеми од претходните трудови.

\textbf{Прово}\cite{Provot1995DeformationCI} користи модел на маса и пружина но додава чекор што прави корекција на рабовите кои се преиздолжени и ги враќа на некоја максимална должина, што ефективно додава ограничувања на должина. Ова е концептуалниот почеток на методите базирани на позиција: наместо да се пресмета само сила на пружина, директно се спроведува дистанцата.

\textbf{Бараф и Виткин}\cite{BaraffDavidWitkin1998} прават имплицитна интеграција на равенките на движење на ткаенини, со што дозволуваат поголеми временски чекори од експлицитните методи. Тука се поставува проблемот на стабилност наспроти цврстина што подоцна PBD целосно го одбегнува.

\textbf{Јакобсен}\cite{jakobsen2001advanced} на неговата GDC презентација за комплексна физика на карактери објаснува за Верле интеграција споена со ограничувања за растојание применети итеративно. Било многу влијателно за игрите, многу независни играчки погони ги имаат изградено своите системи за ткаенини врз основа на оваа рамка пред терминот PBD да постои.

\textbf{Милер, Хејделбергер, Тешнер}\cite{ShapeMatching} го воведуваат методот на спојување на облици спомнат погоре што подоцна се интегрира во PBD системите како еден тип на ограничувања.

\textbf{Милер, Хајделбергер, Хеникс, Ратклиф}\cite{pbd} го воведуваат PBD спомнат погоре што со огромна брзина почнува да се користи во игри и во академијата.

\textbf{Маклин и Милер}\cite{macklin2013position} со \textit{Position Based Fluids}. Додаваат ограничување на густина врз честички, што создава однесување на течност налик на SPH во рамка на PBD. Демонстрираат дека формулацијата базирана на ограничувања генерализира надвор од меки тела и ткаенини.

\textbf{Маклин, Милер, Чентанез}\cite{macklin2016xpbd} со XPBD ја воведуваат формулацијата објаснета погоре и што се користи во овој труд.

После 2017 се намалува брзината на фундаментални промени, понатамошните трудови се фокусираат за додавање на екстензии(тетиви, анисотропни материјали, течности) и инженерски подобрувања(имплементации на GPU, паралелно распоредување на ограничувања, справување со контакти). XPBD останува најкористената формулација за користење во системи во реално време во моментот на пишување на овој труд.

    \begin{equation} \label{BA-cost-function}
      E = \sum_j \rho_j \left( \left\| \pi(\mathbf{P}_c, \mathbf{X}_k) - \mathbf{x}_j \right\|_2^2 \right)
    \end{equation}

\section{Техники на колизии во физички симулации}

Колизиите во физичките симулации имаат своја литература долга децении што е главно независна од литературата за деформации. Спојувањето на двете е главна цел на овој труд.

\subsection{Тврди наспроти меки тела}

Колизија на цврсти тела e добро проучен проблем. За конвексни тела, GJK (Gilbert, Johnson, Keerthi) и EPA(Expanding Polytope Algorithm) пресметуваат растојание и длабочина на пробивање. За произволни триаголнести мрежи, BVH и SAT(Separating Axis Theorem) се стандардот. Контактот помеѓу две цврсти тела се редуцира на мало множество на контактни точки со нормали и длабочини.

Колизијата на меки тела има поголеми побарувања. Геометријата на телото се менува на секоја слика, па структурите за забрзување мора да поддржуваат реизградба или нагодување. Множеството на контакт е огромно: стотици честички може да се во контакт истовремено. Сопствените колизии исто така се од големо значење, телото може да се допира самото себе и широкиот појас на забрзување мора добро да се справи со овој проблем.

\subsection{Дискретна наспроти континуирана детекција}
Дискретните детекции за колизии ги проверуваат позициите на крајот на временскиот чекор, доколку честичката се најде внатре во друго тело, се генерира контакт. Ова е ефтино за проверка и се користи во овој труд. Цената за ова е \textit{тунелирање}, ако честичката се движи доволно брзо може да пројде низ некое доволно тенко тело без да се детектира колизија.

Континуирана детекција на колизии (оттука CCD) го тестира целиот волумен што го опишува честичката при своето движење во временскиот чекор, доколку се сече со некое друго тело се генерираат колизии што би се пропуштиле при дискретната алтернатива. CCD е точно но скапо, најчесто неколу пати поскапо за секоја проверка од дискретна проверка. Модерните игри користат комбинација од дискретни со селективни континуирани детекции за некои објекти

\subsection{Парадигми за реакција на колизија}

Постојат четири главни семејства на реакција при судир, при што секое има свои предности и недостатоци.

\textit{Методите со казни}\cite{hahn1988realistic} го третираат судирот како цврста пружина: силата на реакција е пропорционална на длабочината на продирањето, придвижувајќи ги телата што се судираат настрана. Едноставни се за имплементирање, но цврстината на пружината станува параметар за подесување што се бори против стабилноста на интеграцијата. Методите со казни се во голема мера застарени во современите енџини за меки тела, преживувајќи главно во хаптички симулатори каде што формулацијата со "цврста пружина" се совпаѓа со моделот на хаптичка сила.

\textit{Импулсните методи} (Миртич 1996, Стјуарт 2000) пресметуваат моментална промена на брзината во точката на контакт што спречува меѓусебно продирање, користејќи ги законите за зачувување на контактната механика. Стандарди се за симулатори за цврсти тела (Bullet, ODE, Box2D), каде што равенките за цврсти тела допуштаат импулсни решенија во затворена форма. Помалку се применуваат за меки тела бидејќи претпоставката за "цврстина" се нарушува при секој контакт.

\textit{Одговорот заснован на ограничувања} (Baraff 1989, Bridson 2003) ја третира непродорливоста како ограничување на нееднаквост и решава линеарен комплементарски проблем (LCP) или проектирана Гаус-Сајдел итерација. Точен, добро се справува со натрупување и триење, но е скап и сложен за имплементација.

\textit{Проекцијата на положбата} едноставно ги поместува навлезните честички назад на површината, а потоа ја коригира брзината за да се усогласи. Наједноставниот можен одговор — природен во PBD/XPBD бидејќи интеграторот е заснован на позиции. Ова е техниката што се користи во овој труд.
   
Понова насока е судирите да се формулираат како XPBD ограничувања со сопствени Лагранж множители и усогласеност, со што се интегрира ракувањето со контакт во самиот решавач на ограничувања (Милер\cite{muller2020detailed}). Ова е природен следен чекор надвор од пристапот со проекција на положба што се користи овде, и е дискутиран како иднa работа во Поглавје 9.

\subsection{Пристапи за само-судир}
   
Само-судирот е потежок од судир тело со тело бидејќи телото е подвижен облак од точки, а не фиксен судирач. Тука доминираат три пристапи.

\textit{Просторно хеширање} (Тешнер и сор. 2003\cite{teschner2003optimized}). Има униформна мрежа каде што секоја честичка се хашира во ќелија и се тестираат парови честички во истата или во соседните ќелии. Едноставно, добро прилагодено за деформни тела, користено во оваа теза.

\textit{BVH-ови над геометрија што се деформира}. BVH изграден еднаш во мирување и се прилагодува во секоја слика на деформираната конфигурација. Ефикасно кога деформацијата е умерена, вообичаено во симулатори за ткаенина. Ларсон и Акенин-Молер (2001)\cite{larsson2001collision} воведоа прилагодување од дното нагоре.

\textit{Континуирана само-колизија за ткаенина}\cite{bridson2002robust}. CCD за раб-на-раб и точка-триаголник специјално прилагодено за ткаенина, вклучувајќи зони на удар за справување со кластери што се сечат. Стандард во офлајн симулатори за ткаенина.

Изборот помеѓу нив зависи од тоа дали симулаторот е наменет за дискретно или континуирано детектирање, топологијата на телото (ткаенина наспроти волуменска) и буџетот по слика.



\chapter{Математички и физички основи}

Во овој дел ќе се објаснат математичките и физичките основи кои се користат во останатите глави. 

\section{Њутново движење и временско интегрирање}
Почетната точка е вториот Њутнов закон за секоja честичка:

$$m_i \cdot a_i = F_i$$

каде што $F_i$ е вкупната сила на честичката i. Во овој труд, единствената надворешна сила е гравитацијата, $F_{grav} = m_i \cdot g$, така што интеграторот е одговорен само за напредување на позициите и брзините под дејство на гравитацијата; сите други влијанија (ограничувања на рабови, ограничувања на волумен, судири) се воведуваат како ограничувања во §3.2 и се проектираат за време на решавањето на ограничувањата.

Дискретизацијата во времето создава нумерички интегратор. Три се релевантни.

\subsection{Експлицитен (forward) Ojлер}
Тој е наједноставната шема:

$$v(t + \Delta t) = v(t) + \Delta t \cdot \frac{F(t)}{m}$$
$$x(t + \Delta t) = x(t) + \Delta Δt \cdot (t)$$

Експлицитниот Ојлер е точен до прв ред (грешката се намалува линеарно со бројот на итерации) и многу евтин, но два проблема го прават несоодветен како примарен интегратор за симулатор на честички. Прво, тој не ја зачувува енергијата: вкупната енергија на затворен конзервативен систем расте со текот на долгото симулирање. Едноставен хармониски осцилатор симулиран со eксплицитниот Ojлер покажува видливо спирално однесување во рок од неколку стотици временски чекори. Второ, тој е \textit{нестабилен за тврди системи}: максималниот стабилен временски чекор се скалира како $\Delta t \leq 2/ω$, каде што $\omega$ е највисоката природна фреквенција во системот. Тврдите ограничувања ја зголемуваат $\omega$, принудувајќи на многу мали временски чекори.

\subsection{Симплектички(полуимплицитен) Ојлеров метод}

Се прави мала модификација на експлицитниот Ојлеров метод:

$$v(t + Δt) = v(t) + \Delta t \cdot \frac{F(t)}{m}$$
$$x(t + Δt) = x(t) + \Delta t · v(t + \Delta t) ← ја користи новата брзина$$

Клучната промена е што ажурирањето на позицијата го користи ажурираниот вектор на брзина. Ова го прави методот симплектички интегратор: тој зачувува дискретна аппроксимација на енергијата (и други Хамилтонови инваријанти) на долги временски опсези. Поместувањето на енергијата е ограничено, наместо монотоно.

Симплектичкиот Ојлер е чекорот за предвидување што се користи во PBD и XPBD. Во комбинација со проекција на ограничувањата во чекорот за пресметување на позицијата, тој дава безусловно стабилно однесување дури и за тврди материјали, бидејќи интеграторот не мора директно да ракува со тврдите сили, тој само треба да ги предвиди позициите, кои потоа ги коригира решавачот на ограничувања.

\subsection{Верлеова интеграција}
Трета вообичаена шема е Верлеовата интеграција:

$$x(t + \Delta t) = 2x(t) - x(t - \Delta t) + \Delta t^2 \cdot \frac{F(t)}{m}$$

Со Верле методот не се складира експлицитна брзина; брзината се пресметува од разликата помеѓу последователните позиции. Верле е исто така симплектички и е математички еквивалентен на симплектичкиот Ојлер за системи со константна маса и конзервативни сили. Презентацијата на Јакобсен на GDC во 2001 година го популаризираше Верле во игрите токму затоа што ограничувањата можат да се применат директно на позициите без потреба од следење на брзините, став што PBD го наследи.

Во современите формулации, изборот помеѓу симплектичкиот Ојлер и Верле е во голема мера козметички. Оваа теза ја следи XPBD формулацијата \cite{macklin2016xpbd} во користењето на симплектичкиот Еулер со експлицитно следење на брзината, така што пригушувањето на брзината и триењето се полесни за изразување.


\section{Динамика заснована на ограничувања (Constraint-based dynamics)}
Динамиката заснована на ограничувања е реформација на проблемот со симулација. Наместо да се специфицираат сили и да се врши интегрирање, се специфицира што мора да важи за конфигурацијата: зачувана должина на рабови, зачуван волумен, непенетрирање на честичките и се дозволува решавачот да ја пресмета корекцијата на положбата што ги прави тие тврдења точни.

\subsection{Ограничувања}
Ограничување е скаларна функција $C(X)$ на конфигурацијата на системот. Во оваа теза се појавуваат два вида:
\begin{itemize}
	\item \textit{Ограничувања на еднаквост}: $C(X) = 0$. Примери: должината на работ е еднаква на нејзината должина во мирување; волуменот на тетраедар е еднаков на неговиот волумен во мирување.
	\item \textit{Ограничувања на нееднаквост}: $C(X) \geq 0$. Пример: честичката мора да остане надвор од судирачот, така што \textit{C(X) = означена оддалеченост до површината.}
\end{itemize}


Градиентот на ограничувањето $\nabla C \in \mathbb{R}^{3n}$ ја означува насоката во просторот на конфигурации во која C се менува најбрзо. Повеќето градиенти на ограничувањата се ретки — ограничувањето на работ помеѓу честичките i и j има вредности различни од нула само во блоковите што соодветствуваат на $x_i$ и $x_j$.

\subsection{Лагранжови множители}
Третирајќи го C како холономско ограничување, силата на ограничувањето ја има формата $f = \lambda \nabla C$, каде што $\lambda$ е скаларен Лагранжов множител. Интуицијата е дека ограничувањето го турка системот по насока на $\nabla C$ за количина пропорционална на тоа колку е моментално прекршено. Множителот се одредува со барањето системот да остане на множеството на ограничување $C = 0$.

Со едно ограничување и предвидената положба x* од §3.1, проектираната положба x која го задоволува условот C = 0 (во прв ред) е:

Δx = M⁻¹ ∇C λ λ = −C(x*) / (∇C · M⁻¹ · ∇C)

Ова е линеаризирана проекција на множеството на ограничувања, со тежина од инверзната маса. Истата идеја се генерализира на повеќе ограничувања.

\subsection{Гаус-Сајделова итерација}
Систем од m ограничувања генерира поврзан систем на равенки за множителите. Решавачите од типот PBD не го решаваат поврзаниот систем директно; наместо тоа, тие ја извршуваат Гаус-Сајдел итерацијата, посетувајќи го секое ограничување по ред, независно проектирајќи го и ажурирајќи ја позицијата. Следното ограничување ја гледа ажурираната позиција, така што поврзаноста е фатена имплицитно преку редоследот на евалуација.

Гаус-Сајделова итерација конвергира релативно бавно во споредба со глобалните решавачи (на пр. конјугиран градиент на линеаризираниот систем), но е тривијална за имплементирање, не бара глобално собирање и е лесна за стабилизирање. Тоа е стандарден образец на решач во PBD/XPBD во реално време. Тековната имплементација стандардно користи десет Гаус-Сајделови итерации по потчекор; ова е прикажано како лизгач во графичкиот кориснички интерфејс (GUI).

\section{XPBD формулацијата}

XPBD го дополнува PBD со две идеи: подложност и акумулирани множители. И двете се неопходни за цврстината да биде независна од бројот на итерации. Овој дел ги претставува равенките за ажурирање на XPBD на кои се потпира остатокот од тезата.

\subsection{Усогласеност (compliance)}
Усогласено ограничување е ограничување со конечна крутост. Усогласеноста  $\alpha$ е обратното од крутоста k; $\alpha = 0$ одговара на тврдо ограничување $ (k = \infty)$, а $\alpha > 0$ на меко ограничување. Специфицирањето на материјал во однос на усогласеност е попогодно отколку специфицирањето на крутост, бидејќи $\alpha = 0$ е добро дефинирано, додека $ (k = \infty)$ не е.

Во имплементацијата, и рабовните и волуменските ограничувања имаат изложен параметар на подложност (кој се поставува преку графичкиот кориснички интерфејс) кој се појавува во ажурирањето по итерација, изведено подолу.


\subsection{Терминот alpha}
Кога XPBD се изведува од имплицитна временска интеграција на системот на ограничувања со усогласување, усогласеноста се појавува во равенките за ажурирање поделена со $\Delta t^2$:

$$\tilde{\alpha} = \frac{\alpha}{\Delta t^2}$$

Ова деление одразува дека усогласеноста се компензира со чекорот на времето: ограничување со усогласеност α се однесува поцврсто при помали чекори на времето. Терминот $\tilde{\alpha}$ е она што ја прави цврстината на XPBD независна од итерацијата — тој се појавува експлицитно во ажурирањето по итерација, така што решавачот ја знае "вистинската" целна усогласеност без оглед на тоа колку итерации се потребни.

\subsection{Акумулиран Лагранж мултипликатор}
Во PBD проекцијата во секоја итерација е независна; множителот од претходната итерација се заборава. Во XPBD множителот $\lambda$ се задржува низ итерациите во рамките на потчекор, собирајќи го вкупниот импулс на ограничувањето што е применет. $\lambda$ се ресетира на нула на почетокот на секој потчекор и на секоја итерација се решава само за прирастот $\Delta \lambda$.

Оваа разлика е она што го прави XPBD да конвергира кон одредена цел, наместо да осцилира околу неа.



\subsection{Ажурирањето по итерација}

За едно ограничување C со градиент $\nabla C$ и акумулиран множител $\lambda$, една итерација на XPBD пресметува:

$$\Delta \lambda = -\frac{C(X) - \tilde{\alpha} \lambda}{(\nabla C(x) \cdot M^{-1} \cdot \nabla C(x)^T + \tilde{\alpha})} $$

Читајќи го ажурирањето ред по ред:

Броителот $(C + \tilde{\alpha}\lambda)$  е тековното пречекорување на ограничувањето зголемена за терминот на усогласеност. Како што $\lambda$ расте, зголеменото  се намалува, сигнализирајќи дека ограничувањето се приближува кон задоволување.
Именителот $(\nabla C(x) \cdot M^{-1} \cdot \nabla C(x)^T + \tilde{\alpha})$ е генерализираната инверзна маса (или "ефективна маса") на ограничувањето, плус $\tilde{\alpha}$ Додаденото  $\tilde{\alpha}$ го регулира системот кога геометријата е лошо определена и ја претвора проекцијата од бесконечно цврста во усогласена.
Зголемувањето на положбата $\nabla X $ е во насока на $M^{-1} \nabla C$, скалирано со $\Delta \lambda$. Потешките честички (со помали $w_i$) се движат помалку; полесните честички се движат повеќе.
За $\alpha = 0$, ажурирањето точно се сведува на PBD проекцијата — XPBD е строга генерализација на PBD.

\subsection{Ажурирање на брзината}
Откако сите ограничувања ќе се итерарираат до конвергенција (или ќе се потроши буџетот за итерации), брзините се пресметуваат од промената на позицијата во текот на потчекорот:
$$v = \frac{x - x_{prev}}{\Delta t}$$

Ова имплицитно ажурирање на брзината е карактеристична особина на методите засновани на позиција: силите никогаш не се појавуваат експлицитно во решавачот на ограничувања, туку само позициите и брзините изведени од нив.

\subsection{Зошто функционира цврстината независна од итерацијата}
Интуицијата за цврстината независна од итерацијата: зголеменото прекршување $(C + \tilde{\alpha}\lambda)$ мери колку е системот далеку од целната конфигурација со цврстина, а не од конфигурацијата со строго ограничување. Како што се зголемува $\lambda$, системот се опушта кон рамнотежа со брзина одредена само од $\tilde{\alpha}$, а не од бројот на итерации. Дуплирањето на итерациите го приближува системот до истата цел, наместо да ја поместува самата цел.

\subsection{Потчекори наспроти итерации}
Трудот на Macklin од 2019 година, "Small Steps in Physics Simulation", тврди дека, со даден фиксен вкупен пресметковен буџет, повеќе потчекори со помалку итерации во секој од нив генерално даваат подобри резултати отколку помалку потчекори со многу итерации. Вообичаена конфигурација е 10 потчекори $\times$ 1 итерација, наместо 1 потчекор $\times$ 10 итерации. Имплементацијата ги изложува и двете како параметри; стандардната вредност е 10 $\times$ 10 (цел од 60 fps со удобна маргина).


\chapter{Модел на меко тело}
Кога се импортира некоја мрежа за рендерирање на екран (пример од .obj датотека), во основа се добиваат две листи: листа на 3Д темиња и листа на триаголници кои ги индексираат тие темиња. Тоа е доволно да се исцрта објектот (бидејќи рендерерот ги растеризира триаголниците), но тоа само ја опишува површината на објектот. За рендерирање ова е доволно, но за физички симулации не е.

\subsection{Зошто празна обвивка не може да симулира меко тело}
Како пример можеме да земеме мрежа што репрезентира сфера. Ако наместиме ограничувања за растојание помеѓу соседните темиња (што се рабовите на телото) и со тоа пробаме да симулираме меко тело, добиваме 2 проблеми:

\begin{itemize}
	\item \textbf{Нема задржување на волумен}. Ако се турне едно теме навнатре и системот за ограничувања на телото нема нешто што ќе се спротистави на движењето. Другите темиња остануваат таму каде што се, само рабовите што се директно истегнати или компресирани ја регистрираат промената. Поради ова телото се сплескува како издишан балон кога има било каква сила што турка навнатре
	\item \textbf{Инверзиите се невидливи}. Ако мрежата се извитка низ себеси (надворешна страна пројде низ внатрешна), рабовите на површината остануваат со иста должина и ограничувањата не може да разликуваат дали телото се извртело однадвор навнатре.
\end{itemize}


Обвивката енкодира дека „овој триаголник се наоѓа тука“, но не „овој регион е цврст“. Мекото тело има потреба од второто

\subsection{Тетраедрализација (tetrahedralization)}
Стандардниот начин да се реши проблемот е да се наполни внатрешноста на телото со густо множество на тридимензионални ќелии кои имаат волумен кој може да се измери. На тие волумени може да се дадат ограничувања (да се враќаат кон волуменот во мирување) и телото станува отпорно на компресија навнатре.

Природниот геометриски примитив за ова е \textbf{тетраедарот}, пирамида со 4 темиња што е наједноставното тридимензионално тело со волумен што не е 0 и е добро дефиниран. Тетраедрите имаат неколку особини што ги прават соодветни:
\begin{itemize}
	\item \textbf{Наједноставен примитив}. Три точки дефинираат триаголник, а четири тетраедар. Било што поедноставно е дегенеративно
	\item \textbf{Волумен во затворена форма}. Волуменот е лесен за пресметување со користење на тројниот продукт на три вектори. Нема интеграција и може да се диференцира
	\item \textbf{Едноставни градиенти}. Градиентот на волуменот во однос на секое теме е векторски производ на два спротивни рабови.
	\item \textbf{Има постоечки алатки за пресметка}. Библиотеки како TetGen и CGAL добиваат мрежа од точки како влез и даваат тетраедрализирано тело во секунди
\end{itemize}


\subsection{Делови на мекото тело}
Откако се прави тетраедрализација, мекото тело има четири дела, од кои секој е мотивиран од специфична потреба:

\begin{itemize}
	\item Површински честички (pos, vel, invMass). Телото е облак од точки - нема површина или волумен сè додека ограничувањата не ги декларираат.
	\item Внатрешни честички. Точки кои се додаваат во внатрешноста на телото при тетраедрализација. Рендерерот не ги користи, тие постојат само за да се пресметаат волуменските ограничувања. Сите имаат иста структура како внатрешните точки (позиција, брзина и инверзна маса). 
	\item Рабни ограничувања (едно по уникатен раб на површината). Секое ограничување чува пар на точки со нивната далечина при мирување. Овие ограничувања се спротиставуваат на истегнување и ја дефинираат надворешната форма на објектот.
	\item Волуменски ограничувања. Има едно за секој тетраедар што се добива при тетраедрализација. Секое содржи волумен на тетраедарот при мирување. Овие ограничувата се спротиставуваат на компресија и инверзија.
\end{itemize}

Двете ограничувања имаат различна улога: рабовите ја задржуваат надворешната форма, волуменот ја задржува внатрешната. Рабовите сами прават обвивка што се истегнува, а волумените само прават множество тетраедри без обвивка. Тие заедно даваат тело што може да се деформира.


\section{Репрезентација на честички}
Секоја честичка содржи три информации за состојбата: позиција $x \in \mathbb{R}^3$, брзина $v \in \mathbb{R}^3$, и инверзна маса $w = \frac{1}{m}$. Тие се складираат во \textbf{std::vector} низи во \textbf{SoftBody} struct во include/physics/particle.hpp.

Одлуката да се користи \textit{инверзна} маса е намерна и соодветствува на конвенцијата на PBD/XBPD. Секое ажурирање на ограничувањата дели со сума на инверзни маси, па кеширање на w избегнува делење на секоја итерација. Поважно, w = 0 е природно енкодирање на неподвижна честичка: секоја корекција $\Delta x = w \Delta\lambda \nabla C$ која се применува врз честичката дава нула, па таа останува во иста положба без додавање услови во решавачот. 

Втор пар на низи, \textit{initialPos} и \textit{initialVel} се зачувуваат при градење на објектот. Тие се почетната конфигурација и се користат за ресетирање на симулацијата преку графичкиот интерфејс без одново да се гради објектот.

Топологијата се чува одвоено од состојбата на секоја честичка:
\begin{itemize}
	\item \textbf{tets} - листа од \textbf{array<int, 4>} четворки на индекси, една по тетраедар од волуменската мрежа
	\item \textbf{surfaceIndices} - листа од индексите на темињата од кои се состои мрежата на површинските честички. Се користи од рендерерот за приказ на телото и за ажурирање на позициите на секоја слика
	\item \textbf{numSurfaceParticles} - бројот на уникатни темиња на површината. Се користи од рендерерот за да се исцртаат точниот број на темиња.
\end{itemize}

Распределба на маса: масата по честичка е волуменски нормализирана, при што масата од тетраедрите се префрла на нивните темиња. Секој тетраедар со волумен V придонесува со $\frac{V}{4}$ од волуменот на телото во секое од своите четири темиња, па акумулираната вредност за честичката е вкупниот волумен на материјал што таа го претставува. Аргументот за маса во конструкторот е посакуваната вкупна маса на телото; еднакво распределена густина $\rho = \frac{mass}{V_{total}}$ го претвора акумулираниот волумен во маса по честичка, давајќи $w_i = \frac{1}{(\rho \cdot V_i)}$. Ова е пофизичко од еднакви маси, честичките во поголемите или погусти региони на тетраедрите имаат поголема инерција, додека честичките во тенките региони се полесни, и ја одржува вкупната маса на телото независнo од скалата на мрежата. Волуменот на секој тетраедар го реупотребува скаларниот троен проивод кој подоцна се користи од страна на ограничувањето за волумен (§4.4). Честичките кои не припаѓаат на ниту еден тетраедар, и дегенерирана мрежа со вкупен волумен нула, се враќаат на w = 0 (прицврстени), во согласност со конвенцијата за прицврстување со инверзна маса погоре.


\section{Тетраедрализација}
\subsection{TetGen}

Волуменската мрежа во симулатор е произведена од TetGen \cite{tetgen}, генератор на тетраедрални мрежи базиран на Делоне метод. TetGen на влез зема PLC (piecewise linear complex) и како резултат дава ограничена Делоне тетраедрализација: се додаваат внатрешни Штајнер точки по потреба и волуменот помеѓу внатрешните страни се пополнува со тетраедри што го задоволуваат Делоне условот на празна впишана сфера.

Во симулаторот TetGen е интегриран на следниот начин:
\begin{enumerate}
	\item Влезните темиња се запишуваат во \textbf{in.pointlist} како рамна низа
	\item Површинските триаголници се запишуваат во \textbf{in.facetlist} како страни од еден полигон
	\item Се конфигурира \textbf{tetgenbehavior} со знамето \textbf{pq2.0Y}
	\item \textbf{tetrahedraliza(\&b, \&in, \&out)} го извршува алгоритмот
	\item Излезот \textbf{out.pointlist, out.tetrahedronlist, и out.trifacelist} се користат од низите \textbf{pos, tets, surfaceIndices} од симулаторот
\end{enumerate}

Стрингот \textbf{pq2.0Y} што се користи за конфигурација се состои од:
\begin{itemize}
	\item \textbf{p} - \textit{PLC мод}. Му кажува на TetGen дека влезот не е облак од точки туку полиедар опишан од темиња и страни, па само внатрешноста треба да се тетраедрализира
	\item \textbf{q1.4} - \textit{рафинирање на квалитет} со сооднос на радиус-раб од 1.4. Овој сооднос е опишаниот радиус на тетраедарот поделен со должината на неговиот најкраток раб. Пониски вредности произведуваат повеќе тетраедри со цена на повеќе Штајнер точки и повеќе тетраедри на излез. 
	\item \textbf{Y} - \textit{задржување на влезната граница}. Без ова знаме, TetGen може да внесува Штајнер точки на површинските триаголници, што ќе ја наруши претпоставката дека влезните индекси на темиња се мапираат со првите позиции во низата \textbf{pos}.
\end{itemize}

\subsection{Ограничувања}
TetGen зависи од точен влез. Ако влезната мрежа не е затворена, има пресеци сама со себе, или ако има погрешна ориентација на страните, ќе се добие мрежа со дупки, грешки при извршување или тивок неуспех каде површината е точна но внатрешноста остане празна. Доколку се добие некој погрешен резултат со нов објект додаден во сцената, погрешниот тип на објект е прво нешто што треба да се провери.

\section{Ограничувања на рабови}
Ограничувањата на рабови се „надворешната“ половина на мекото тело. Во овој дел се дава формална дефиниција, градиентот кој се решава и кратка дискусија за однесувањето и подесување.

\subsection{Дедупликација}
Површинските триаголници имаат заеднички рабови, па наивно креирање на едно ограничување по раб на триаголник би креирално дупликати на рабовите. \textbf{Solver::createConstraintFromIndices} прави дедупликација со внесување на парови на индекси од темиња и скокање на веќе постоечките парови.
 
\subsection{Ограничување и градиент}

Функцијата на ограничувањето е отстапувањето од должината при мирување:

$$C(x_1, x_2) = \lVert x_1 - x_2 \rVert - I_0$$

Должината при мирување $I_0$ се пресметува при конструирање на телото од иницијалните позиции на честичките, па „обликот при мирување“ на телото е геометријата што се вчитува на почеток на симулацијата.

Градиентот е единечен вектор со насока паралелна на работ:
$$\nabla_{x_1} C = \hat{n}, \nabla_{x_2}C = -\hat{n}$$
$$ \hat{n} = \frac{(x_1 - x_2)}{\lVert x_1 - x_2 \rVert}$$

Градиентите имаат единечна должина, што прави броителот во XPBD да се скрати до $w_1 + w_2 + \tilde{\alpha}$, што е ефтино за пресметување.

\subsection{Чекор во решавач}

Во \textbf{Solver::solveEdgeConstraint} се имплементира ажурирањето на позициите и ограничувањата:
\begin{enumerate}
	\item Читање на позиции и инверзни маси. Излез ако двете честички имаат бесконечни маси
	\item Пресметка на $I = \lVert x_1 - x_2 \rVert$. Излез ако $I < 10^{-6}$ (Дегенериран раб со недефиниран градиент)
	\item Нормализирање $\hat{n} = \frac{x_1 - x_2}{I}$
	\item $C = I - I_0$
	\item $\tilde{\alpha} = \frac{\alpha}{\Delta t^2}$
	\item $\Delta \lambda = \frac{-C-\tilde{\alpha}\lambda}{(w_1 + w_2 + \tilde{\alpha})}$
	\item $\Delta x_1 = w_1 \Delta \lambda \hat{n}; \Delta x_2 = -w_2 \Delta \lambda \hat{n}$
	\item $\lambda \leftarrow \lambda + \Delta \lambda$
\end{enumerate}

Проверките во чекор еден и два за ран излез се неопходни. Без нив, две честички кои се неподвижни ќе предизвикаат делење 0/0, а целосно дегенериран раб ќе произведе недефинирана насока

\subsection{Ограничување на површински рабови}

Само рабовите од надворешната обвивка на телото имаат ограничувања. Внатрешните кои ги поврзуваат Штајнер точките до нивните соседи немаат ограничување на раб, тие се контролираат само од волуменските ограничувања. Друга опција е да се додадат ограничувања за секој раб во театраедрализираната мрежа, но тоа ќе произведе поцврсто тело што поагресивно конвергира при истегнување, но истовремено ќе се додадат неколкупати повеќе ограничувања.

\subsection{Однесување и нагодување}
Со $\alpha = 0 $ (или многу мала вредност), ограничувањата на раб се однесуваат скоро како цврсти ограничувања на растојание, при што телото се спротиставува на истегнување. Со поголемо $\alpha$, рабовите стануваат како гума, телото повеќе се сплескува под гравитација, и осцилациите од висока фреквенција брзо се намалуваат. Бидејќи усогласеноста влага во именителот, поголемо $\alpha$ го прави системот постабилен.


\section{Волуменски ограничувања}
Волуменски ограничувања се „внатрешната“ половина на мекото тело, тие предизвикуваат тоа да се спротиставува на компресија и инверзија, што рабовните ограничувања не можат да го направат. Во овој дел се опишува функцијата на ограничување, формулите за  градиенти, чекорот во решавачот.

\subsection{Скаларен троен производ}
Волуменот на тетраедар со темиња $\mathbf{x_1, x_2, x_3, x_4}$ е:
$$V = \frac{1}{6}\mathbf{(x_2 - x_1)}\cdot((\mathbf{x_3 - x_1})\times\mathbf{(x_4 - x_1)})$$

Во имплементацијата се зачувува $6V_0$ за да се прескокне делењето со $\frac{1}{6}$ во секоја итерација. Конкретно, \textbf{restVolume} во кодот е тројниот продукт на трите вектори на рабовите при конструкција на објектот.

Функцијата на ограничување е 
$$ C = 6V - 6V_0$$

при што C > 0 ако тетраедарот е поголем од почетниот волумен, C < 0 ако е копмресиран или инвертиран.

\subsection{Градиенти}
Градиентите се векторските производи на спротивните вектори од рабовите на секое теме, односно:
$$\nabla_1 C = \mathbf{(x_4 - x_2) \times (x_3 - x_2)}  $$
$$\nabla_2 C = \mathbf{(x_3 - x_1) \times (x_4 - x_1)}  $$
$$\nabla_3 C = \mathbf{(x_4 - x_1) \times (x_2 - x_1)}  $$
$$\nabla_4 C = \mathbf{(x_2 - x_1) \times (x_3 - x_1)}  $$
Геометриски, $\nabla_i C$ покажува надвор од страната која е спроти темето \textit{i}, со должина еднаква на двапати од плоштината на таа страна. Физички, влечење на темето \textit{i} нанадвор во таа насока го зголемува волуменот.

\subsection{Чекор во решавачот}
\textbf{solveVolumeConstraints} имплементира:
\begin{enumerate}
	\item Читање на сите позиции и инверзни маси, излез ако некое теме е неподвижно
	\item Пресметка на градиентите од погоре
	\item Добивање на моментален 6V директно од секој од нив: $6 = \nabla_4 \cdot(\mathbf{x_4 - x_1}) $(избегнува повторно пресметување на тројниот продукт)
	\item $C = 6V - 6V_0$
	\item $\tilde{\alpha} = \frac{\alpha}{\Delta t^2}$
	\item $D =  w_1 \lVert \nabla_1 \rVert^2 + w_2\lVert \nabla_2\rVert^2 + w_3\lVert\nabla_3\rVert^2 + w_4\lVert\nabla_4\rVert^2 + \tilde{\alpha}$
	\item Излез ако $D<10^{-8}$ (дегенеративен тетраедар)
	\item $\Delta \lambda = \frac{(-C -\tilde{\alpha}\lambda)}{D}$
	\item $\Delta x_i = w_i \Delta \lambda \nabla_i $ за $i = 1..4$
	\item $\lambda \leftarrow \lambda + \Delta \lambda$
\end{enumerate}

\subsection{Инверзија и дегенерација}
Тетраедар чии темиња се копланарни има нулти волумен и нулта големина на градиентот, при што именителот е само  $\tilde{\alpha}$, па ажурирањето дегенерира. Тетраедар што прошол во состојба каде што има негативен волумен (односно е инвертиран)  сѐ уште има добро дефинирани градиенти, а ограничувањето го повлекува назад во валидна состојба. Во практика инверзијата е ретка кога се користат ограничувања на рабови и волумени при нормални временски чекори, кога се случува обично е знак дека временските чекори се предолги или усогласувањето е превисоко.

\subsection{Сооднос од 0.1 за усогласување}
Во конструктурот на \textbf{Solver::Solver}, волуменското усогласување е наместено да е 0.1$\times$ рабното усогласување. Ова е хардкодирано, можеби ако се даде како посебен параметар што може дасе прилагоди е идно подобрување......


\section{Циклусот на XPBD решавачот}
\textbf{Solver::update(dt)} е главната функција на целата симулација. Секој повик го зголемува времето за dt(што е времето поминато од една до друга слика) и е поделен на \textbf{substeps} број на подчекори. Внатре во секој подчекор се извршува процес од пет фази:

\begin{algorithm}[H]
	\normalsize                         % <-- shrinks everything uniformly
	\SetCommentSty{scriptsize}      % <-- comments one size smaller still
	\DontPrintSemicolon
	\caption{\textsc{Solver::update}$(dt)$}
	
	\KwIn{Time step $dt$, substeps $S$, solver iterations $K$}
	
	$h \leftarrow dt / S$\;
	$e_r \leftarrow 0.3$\tcp*{restitution}
	$\mu   \leftarrow 0.7$\tcp*{friction}
	
	\For{$s = 1$ \KwTo $S$}{
		
		\tcp{1.\ Prediction (symplectic Euler)}
		\For{\textnormal{each particle} $i$}{
			$\mathbf{v}_i      \leftarrow \mathbf{v}_i + h\,\mathbf{g}$\;
			$\mathbf{x}^{\mathrm{prev}}_i \leftarrow \mathbf{x}_i$\;
			$\mathbf{x}_i      \leftarrow \mathbf{x}_i + h\,\mathbf{v}_i$\;
		}
		
		\tcp{2.\ Spatial-hash refresh}
		$\mathcal{H}.\textsc{clear}()$\;
		\For{\textnormal{each particle} $i$}{
			$\mathcal{H}.\textsc{insert}(i,\,\mathbf{x}_i)$\;
		}
		
		\tcp{3.\ XPBD setup — reset accumulated multipliers}
		\For{\textnormal{each constraint} $c$}{
			$\lambda_c \leftarrow 0$\;
		}
		
		\tcp{4.\ Constraint solve (Gauss–Seidel)}
		\For{$k = 1$ \KwTo $K$}{
			\For{\textnormal{each edge constraint} $c$}{\textsc{SolveEdge}$(c,\,h)$\;}
			\For{\textnormal{each volume constraint} $c$}{\textsc{SolveVolume}$(c,\,h)$\;}
		}
		
		\tcp{5a.\ Collision detection + position correction}
		\For{\textnormal{each particle} $i$}{
			\lIf{$x_{i,y} < 0$}{$x_{i,y} \leftarrow 0$}
			\For{\textnormal{each collider} $\mathcal{C}$}{
				$(d,\,\mathbf{n}) \leftarrow \mathcal{C}.\textsc{SignedDistance}(\mathbf{x}_i)$\;
				\lIf{$d < 0$}{$\mathbf{x}_i \leftarrow \mathbf{x}_i - d\,\mathbf{n}$}
			}
		}
		
		\tcp{5b.\ Velocity update (purely positional)}
		\For{\textnormal{each particle} $i$}{
			$\mathbf{v}_i \leftarrow \bigl(\mathbf{x}_i - \mathbf{x}^{\mathrm{prev}}_i\bigr) / h$\;
		}
		
		\tcp{5c.\ Velocity correction at contacts (restitution + friction)}
		\For{\textnormal{each particle} $i$}{
			\If{$x_{i,y} \leq 0$}{
				$v_{i,y} \leftarrow |v_{i,y}|\cdot e_r$\;
				$v_{i,x} \leftarrow v_{i,x}\cdot\mu$\;
				$v_{i,z} \leftarrow v_{i,z}\cdot\mu$\;
			}
			\For{\textnormal{each collider} $\mathcal{C}$}{
				$(d,\,\mathbf{n}) \leftarrow \mathcal{C}.\textsc{SignedDistance}(\mathbf{x}_i)$\;
				\If{$d < d_{\mathrm{contact}}$}{
					$\mathbf{v}_n \leftarrow (\mathbf{v}_i \cdot \mathbf{n})\,\mathbf{n}$\;
					\If{$\mathbf{v}_i \cdot \mathbf{n} < 0$}{
						$\mathbf{v}_t \leftarrow \mathbf{v}_i - \mathbf{v}_n$\;
						$\mathbf{v}_i \leftarrow -e_r\,\mathbf{v}_n + \mu\,\mathbf{v}_t$\;
					}
				}
			}
		}
	}
\end{algorithm}


\chapter{Техники за детекција на колизии}


   